\chapter{Úvod}
Testovanie je jedna z najdôležitejších fáz vývoja softvéru, pretože dôsledky nekvalitného softvéru sú nákladné. Nekvalitný softvér spôsobil v roku 2022 v USA náklady odhadované na skoro dva a pol bilióna amerických dolárov[tu pridaj citacu krasner 2022]. Jedným z prístupov ku testovaniu je Data-Driven testing (DDT), testovanie založené na princípe oddelenia testovacích dát od testovacej logiky.

Efektívne testovanie stojí na kvalite testovacích dát, hlavne preto, že pri nedostatočne domyslených dátach môže dôjsť k neodhaleným chybám, ale aj k falošnému pocitu dobre otestovaného produktu, čo môže spôsobiť stratu časových a finančných zdrojov.
Vytváranie testovacích dát je časovo náročné a vyžaduje znalosť schémy, ale aj charakteru reálnych dát.

Pre uľahčenie práce existujú nástroje na generovanie náhodných testovacích dát. Jedným z nich je Dataster, nástroj, ktorý generuje štruktúrované testovacie dáta pre relačné databázy na základe konfiguračného receptu. Ručné vytvorenie receptu vo forme YAML vyžaduje znalosť dokumentácie programu Dataster, ako aj znalosť schémy databázy, vzťahov medzi tabuľkami a charakteru uložených hodnôt. Náročnosť tohto procesu rastie so zložitosťou databázy. Pre používateľa s nízkou znalosťou konkrétnej databázy predstavuje tvorba receptu netriviálnu manuálnu prácu.

Potenciálnym riešením je automatizácia tejto úlohy s pomocou veľkých jazykových modelov (LLM), ktoré majú schopnosť interpretovať štruktúrované dáta a transformovať ich do iných formátov. Táto práca využíva lokálne LLM, keďže vstupom sú reálne dáta z databázy a lokálny beh zabezpečuje nielen súkromie dát, ale aj nízke náklady.

Cieľom tejto práce je navrhnúť a implementovať nástroj, ktorý sa pripojí k databáze, extrahuje jej schému a vzorku uložených dát a pomocou lokálnych LLM automaticky vygeneruje recept pre nástroj Dataster. Ďalším cieľom je zhodnotiť kvalitu generovaných receptov naprieč rôznymi modelmi a konfiguráciami. Kvalita sa meria od štrukturálnej správnosti vygenerovaného receptu (YAML) cez akceptovanie nástrojom Dataster až po vykonateľnosť výsledného SQL kódu. Vyhodnotenie prebieha na sade testovacích databáz rôznej zložitosti.

Kapitola 2 približuje definíciu testovania a popis data-driven testing. Kapitola 3 definuje LLM, ich využitie pri generovaní dát, predstavuje jednotlivé LLM modely a skúmanie alternatív. Kapitola 4 približuje nástroj Dataster, jeho využitie, limitácie, popis konfiguračného receptu a YAML súborov. Kapitola 5 predstavuje vysokoúrovňový návrh programu. Kapitola 6 sa zaoberá koncovou implementáciou. Kapitola 7 vysvetľuje proces experimentov, sadu testovacích databáz, úrovne validácie vygenerovaných receptov a porovnávané modely a konfigurácie. Kapitola 8 prezentuje a porovnáva výsledky experimentov. Kapitola 9 objasňuje limitácie programu a možné rozšírenia v budúcej práci. Na záver kapitola 10 zhŕňa dosiahnuté výsledky a porovnáva ich s cieľmi práce.



\chapter{Testovanie softvéru}

\section{Význam testovania}

Testovanie systému je možné chápať ako proces vyhodnotenia systému sledovaním jeho behu. Jeho cieľom je minimalizovať počet chýb a overiť  či software skutočne spĺňa technické požiadavky pri jeho návrhu a vývoji a či uspokojuje užívateľské požiadavky efektívne a účinne. Je to nesmierne dôležitý proces pri vyvoji softveru.


\section{Dôsledky ndôkladného testovania}
Každý softvér je v nejakej fáze vývoja chybný. Tieto chyby možu mať za následok závažné straty, ktoré začínajú od finančných strát až po straty na ľudských životoch. Existuje množstvo zaznamenaných prípadov kde nedostatočné a nedôkladné testovanie viedlo až ku katastrofickým následkom.
Jeden z týchto prípadov je explodovanie rakety len 37 sekund po jej starte spôsobené neošetrenou vynimkou pri konverzii 64-bitového float čísla na 16-bitové signed integer v navigačnom systéme, ktorá viedla k strate rakety a nákladom vo výške až 500 miliónov eur\cite{Smrcka2024}

Ďalším príkladom je Therac-25, prístroj určený na rádioterapiu, ktorý bol používaný v rokoch 1985 až 1987.  Jedenásť prístrojov bolo nasadených v USA a Kanade a v roku 1987 boli odvolané. Chybný softvér spôsobil niekolko prípadov predávkovania radiáciou a aspoň 6 prípadov skončilo so závažným zranením až smrťou. Podobné chyby boli nájdené aj iných Therac modeloch, ale Therac-20 obsahovala bezpečnostný blok a tým pádom k zraneniam nedošlo.\cite{Leveson1993}

Podobných prípadov je ale viac. Poukazujú na dôležitosť testovania, ktoré je nevyhnuteľné pre zabezpečenie kvalitného, bezpečného a spoľahlivého softvéru.


\subsection{Úrovne testovania}\cite{Smrcka2024}

Testovanie prebieha na rôznych úrovniach pričom každá overuje iné aspekty systému:

\begin{itemize}
\item \textbf{Unit testing (jednotkové testovanie)} -- testovanie najmenších jednotiek kódu ako funkcie a metódy.
\item \textbf{Module testing (testovanie modulov)} -- testovanie kolekcie súvisiacich jednotiek, ktoré riešia spoločný problém. Testovaním sa hľadajú chyby v komunikácii medzi jednotkami.
\item \textbf{Integration testing (integračné testovanie)} -- Overenie rozhrania a komunikacia medzi modulmi.
\item \textbf{System testing (testovanie systému)} -- Kontroluje softvér ako celok. Hľadajú sa chyby v návrhu a špecifikácií.
\item \textbf{Acceptance testing  (akceptačné testovanie)} -- Overenie či softvér spĺňa požiadavky zákazníka a či spĺňa svoj primárny účel. 
\end{itemize}

\subsection{Automatizované testovanie}
Pri väčších projektoch sa manuálne testovanie stáva neefektívne z, nielen, časového hľadiska ale aj ekonomického. Používa sa preto automatizovane testovanie pomocou využitia rôznych frameworkov a nástrojov na automatické vykonanie testov. Prináša mnoho výhod ako:
\begin{itemize}
    \item rýchlejšie vykonanie testov
    \item možnosť častého opakovania testov
    \item vyššie pokrytie kódu
\end{itemize}
Automatizované testovanie sa stáva nevyhnuteľnou súčasťou vývoja. Jedným z pristupov  testovania, s ktorým sa zameriavam aj v tejto práci je Data-Driven Testing (DDT).

\section{Data-Driven Testing}
Spôsob testovania, ktorý oddeľuje testovaciu logiku od testovacích dát.\cite{Laukkanen2006} Základnou myšlienkou je, že rovnaký testovací skript možno spustiť s rôznymi sadami vstupných dát, čím sa  zvyšuje pokrytie testov bez nutnosti duplikovania kódu.\cite{Myers2011}

\subsection{Hlavné charakteristiky}

\begin{itemize}
\item \textbf{Oddelenie dát od logiky} -- testovacie dáta sú uložené v externých súboroch ako napríklad Excel, CSV, YAML, nie v kóde testu.
\item \textbf{Znovupoužiteľnosť} -- rovnaké testy možno spustiť s rôznymi sadami dát
\item \textbf{Čitateľnosť} --  testovacie dáta musia byť ľahko čitateľné aj osobou bez znalosti programovania \cite{Laukkanen2006}

\item \textbf{Jednoduchá údržba} -- možné pridanie nového testu bez zmeny kódu
\end{itemize}

\subsection{Výhody a nevýhody Data-Driven Testing}

Data-Driven testing prináša mnoho výhod ako napríklad zvýšenie pokrytia testov, znovupoužiteľnosť a jednoduchá údržba.

Má však aj svoje nevýhody. Je potrebné dátové sady neustále udržiavať a príprava kvalitných testovacích dát môže byť náročná. Tento typ testovanie je na nich ale závislý.


\section{Existujúce nástroje}

\subsection{Dataster}
Jeden z komponentov Testos platformy, ktorá podporuje automatizáciu testovania softvéru. Slúži na generovanie dát pomocou YAML konfigurácie. Jeho vysledkom je generovanie SQL INSERT príkazov do databázových tabuliek.

Manuálne vytvorenie YAML konfigurácie je nielen časovo náročné ale vyžaduje aj schopnosť manuálnej analýzy databázovej schémy pre identifikáciu dátových typov a definovanie vzťahov medzi tabuľkami. Tento proces vyžaduje hodiny až dni práce a so stúpajúcim množstvom tabuliek je to prakticky nemožné.  Zjednodusenim tohto problemu sa zaobera tato práca.

\section{Veľké jazykové modely (Large Language Models, LLM)}
Large Language Mode, skrátene LLM, je typ umelej inteligencie ktorý dokáže rozpoznať a generovať text ako aj vykonávať iné úlohy. Tieto modely sú trénované na nespočetnom množstve dát. Medzi najznámejšie LLM patrí GPT (OpenAi), Gemini (Google), Llama (Meta).

\subsection{Trénovanie LLM}
Vo fáze trénovania sa LLM učia rozpoznávať štýl a vzor jazyka spracovaním veľkého množstva textu. Učia sa gramatiku, fakty, logicke vztahy a kontextove suvislosti.
Kvalita týchto dát vo veľkej miere ovplyvňuje samotný výkon LLM modelu.\cite{Zhao2024}

Trénovanie LLM možno rozdeliť do troch krokov:\cite{Zhao2024}
\textbf{zber a spracovanie dát:} zoskupenie veľkého množstva dát ako sú napríklad knihy, webové stránky a iné.

\textbf{pre-training:} určenie architektúry modelu a a pred trénovacích úloh s využitím vhodných trenovacich algoritmov


\textbf{fine-tuning a alignment:} doladenie modelu pre konkrétne úlohy

\subsection{Lokálne LLM - Ollama}

Táto práca využíva lokálne LLM pre ich schopnosť analyzovať komplexné dátové štruktúry a pomocou nich následne vygenerovať štruktúrovaný výstup. Toto je využité pre automaticke generovanie Dataster konfigurácie na zaklade databázovej schémy.

Ako lokálne LLM som si vybrala pre Ollama v navrhovanom rieseni pre jeho jednoduchosť inštalácie a open-source a aktívnu komunitu.Ďalším dôvodom je mnohé databázové schémy obsahuju citlivé informácie a tak použitím lokálneho LLM je možné zabrániť prístupu nezavislych stran.

\chapter{Návrh systému}

\section{Všeobecný prehľad riešenia}

Navrhovaný nástroj slúží pre vytvorenie YAML konfigurácie pre generátor testovacích dát Dataster. Jeho cieľom je nahradit manuálne písanie YAML konfigurácie, znížiť časovú záťaž procesu analýzý  databázy a pomocou veľkého jazykového modelu vytvoriť konfiguráciu automaticky.

Kľúčové je ale využitie LLM vo dvoch fázach procesu:
\begin{itemize}
    \item \textbf{Generovanie SQL požiadavkov} --   LLM na základe šablónových SQL požiadavkov vytvorí konkrétne SQL požiadavky pre získanie štruktúry databázy a vzorových dát.
    \item \textbf{Generovanie YAML konfigurácie} -- LLM na základe získaných informácií vytvorí YAML konfiguráciu pre Dataster
\end{itemize}


\section{Architektúra systému}

Systém sa skladá z piatich hlavných komponentov, ktorých vzájomné vzťahy znázorňuje obrázok \ref{fig:lefttoright}.

\begin{figure}[h]
\centering \includegraphics[width=0.9\textwidth]{obrazky-figures/lefttoright.pdf}
\caption{Architektúra navrhovaného systému}
\label{fig:lefttoright}
\end{figure}

\subsection{Popis komponentov}

\subsubsection{DB Connector}

Zabezpečuje pripojenie k relačnej databáze.

\subsubsection{Schema Extractor}

Získava štruktúru databázy pomocou SQL požiadavkov generovaných LLM.


\subsubsection{Data Sampler}

Získava reprezentatívne vzorky dát pomocou SQL požiadavkov generovaných LLM.


\subsubsection{LLM Processor}

Komunikuje s lokálnym LLM modelom (Ollama) a má dve hlavné funkcie:
\begin{enumerate}
    \item \textbf{ Generovanie SQL požiadavok}

    \item \textbf{Generovanie YAML konfigurácie}
\end{enumerate}


\subsubsection{YAML Generator}

Validuje a ukladá vygenerovaný YAML súbor.


\section{Tok dát v systéme}

Celý proces generovania YAML konfigurácie prebieha v nasledujúcich krokoch (viď obrázok \ref{fig:chart}):

\begin{figure}[h]
\centering
\includegraphics[width=0.9\textwidth]{obrazky-figures/chart.pdf}
\caption{Diagram postupnosti operácií v systéme}
\label{fig:chart}
\end{figure}

\begin{enumerate}
\item \textbf{Pripojenie:} Používateľ zadá prihlasovacie údaje, DB Connector vytvorí pripojenie k databáze.

\item \textbf{Získanie zoznamu tabuliek:}
\begin{itemize}
\item Schema Extractor vytvorí šablónovú požiadavku
\item LLM Processor vygeneruje SQL požiadavky
\item Schema Extractor vykoná SQL požiadavku a získa zoznam tabuliek
\end{itemize}

\item \textbf{Získanie štruktúry tabuliek:}
\begin{itemize}
\item Pre každú tabuľku Schema Extractor vytvorí šablónovú požiadavku
\item LLM Processor vygeneruje SQL požiadavky 
\item Schema Extractor vykoná SQL požiadavky a získa informácie o stĺpcoch, typoch a kľúčoch
\end{itemize}

\item \textbf{Vzorkovanie dát:}
\begin{itemize}
\item Pre každú tabuľku Data Sampler vytvorí šablónovú požiadavku
\item LLM Processor vygeneruje SQL požiadavky
\item Data Sampler vykoná SQL požiadavky a získa vzorové riadky pre ilustráciu skutočných hodnôt
\end{itemize}

\item \textbf{Vytvorenie promptu pre YAML:} LLM Processor zozbiera všetky informácie (schému databázy + vzorové dáta) a skonštruuje komplexný prompt obsahujuci popis databázy.

\item \textbf{Generovanie YAML:} Ollama LLM vygeneruje YAML konfiguráciu na základe poskytnutých informácií.

\item \textbf{Validácia:} YAML Generator overí správnosť vygenerovaného YAML.

\item \textbf{Uloženie:} Validovaná konfigurácia sa uloží do súboru \texttt{dataster\_config.yaml}, ktorý je pripravený na použitie s Dataster.
\end{enumerate}


\section{Použité technológie }

\textbf{Programovací jazyk:} Python 3.10+


\textbf{Knižnice:}
\begin{itemize}
\item \texttt{psycopg2} alebo \texttt{mysql-connector-python} -- databázové konektory
\item \texttt{ollama} -- Python SDK pre Ollama
\item \texttt{pyyaml} -- spracovanie YAML súborov
\item \texttt{argparse} -- CLI rozhranie
\end{itemize}

\textbf{Externé nástroje:}
\begin{itemize}
\item \textbf{Ollama} -- lokálny LLM runtime
\item \textbf{Dataster} -- generátor testovacích dát
\end{itemize}

\textbf{Podporované databázy:}
\begin{itemize}
\item PostgreSQL 12+
\item MySQL 8.0+
\end{itemize}


\chapter{Záver}

Pre semestrálny projekt bol navrhnutý program na automatické generovanie YAML konfigurácií pre nástroj Dataster s využitím lokálnych veľkých jazykových modelov.

V rámci projektu boli dosiahnuté požiadavky ako naštudovanie Data-Driven testing, testovanie založené na dátach a zoznámila som sa s generovaním dát pomocou LLM. Následne som analyzovala požiadavky pre automatické generovanie náhodných štruktúrovaných dát určených ako testovacie vstupy. A s týmito znalosťami som navrhla nástroj, ktorý využíva lokálne LLM pre generovanie týchto testovacích dát.

Navrhnutý systém využíva LLM v dvoch kľúčových fázach: najprv na generovanie šablónových SQL požiadavkov pre analýzu databázovej schémy a získanie vzorových dát, následne na vytvorenie YAML konfigurácie na základe získaných informácií. Tento dvojfázový prístup umožňuje plnú automatizáciu procesu, ktorý by inak vyžadoval manuálnu prácu v rozsahu hodín až dní.

\section{Ďalšie kroky}

Pri pokračovaní bakalárskej práce bude návrh implementovaný ako nástroj s využitím zvoleného lokálneho LLM podľa navrhnutej architektúry. Následne testovanie správnosti s pomocou rôznych databáz s rôznou komplexnosťou a veľkosťou a s týmito znalosťami vyhodnotenie kvality riešenia pre rôzne LLM modely.